\documentclass[11pt]{amsart}
\usepackage{amsmath,amssymb,amsthm,hyperref}
\usepackage[margin=1in]{geometry}

\title[Darmon--Merel for $a^4+b^4=c^{13}$]{A formalization of the Darmon--Merel method for $a^4+b^4=c^{13}$}
\author{David Fox}
\thanks{Source: \url{https://github.com/DavidFox998/beal-conjecture}, branch \texttt{mathlib/remove-ribet-axiom}, tag \texttt{v39-prototype-4413}.
Lean~4.34.1, Init only, no community Mathlib.
This note formalizes that method for one signature.
It is not a proof of $a^4+b^4=c^{13}$, and it is not a proof of the Beal conjecture.
It may be published on its own, at the editors' discretion.}
\date{}

\begin{document}
\begin{abstract}
We record a Lean~4 formalization of the Darmon--Merel method for the equation $a^4+b^4=c^{13}$.
The method attaches a Frey $\mathbb Q$-curve over $\mathbb Q(i)$ and compares a lowered newform at level $32$ or $64$ with a trace.
The special-fiber blow-up is counted: seven components when $a$ is odd, and nine when $a$ is even.
Pacetti's Artin exponent is then a theorem of the two propositions that identify those counts with the N\'eron special fiber.
The mod-$13$ residual and the $\mathbb Q$-curve extension are present as definitions and statements.
The nonexistence statement for coprime $a,b$ both nonzero remains conditional on those two fiber propositions, on the absence of a newform of level $32$ or $64$, on an opaque irreducibility hypothesis, and on two named axioms.
A further trace hypothesis over $\mathbb F_{169}$ implies the newform proposition and is not assumed by the capstone.
There is no \texttt{sorry}.
This is not a proof of $a^4+b^4=c^{13}$, and it is not a proof of the Beal conjecture.
\end{abstract}
\maketitle

\section{What is formalized}
The equation is the Darmon--Merel setting of Pacetti's equation $x^4+dy^2=z^p$ at $d=1$ and $p=13$, with $x=a$, $y=b^2$, and $z=c$.
The curve is
\[
y^2 = X^3 + 4a X^2 + 2\bigl(a^2 + b^2\sqrt{-1}\bigr) X,
\]
with discriminant $\Delta = 2^9(a^2+b^2\sqrt{-1})\,c^{13}$.
The checked library is the branch and tag named in the thanks.
\texttt{lake build} and \texttt{lake build Mathlib} are the two build targets.
The printed axioms of the capstone \texttt{no\_solution\_4413\_of\_frontier} are \texttt{frey\_qcurve\_modular}, \texttt{ribet\_level\_lowering\_4413}, \texttt{propext}, \texttt{Classical.choice}, and \texttt{Quot.sound}.
There is no \texttt{sorry} in the library.

The pairs $(1,0,1)$ and $(0,1,1)$ satisfy the equation, with $\gcd=1$ and $c>0$.
They are excluded by the quantifiers $a\neq 0$ and $b\neq 0$.
The pair $(8,8,2)$ fails coprimality.
There is no theorem that $a$ and $b$ cannot both be nonzero.

\section{The two component counts}
For odd $a$, the charts in \texttt{Mathlib/Modular/BlowupOddA.lean} give
\[
\texttt{OddAComponent.all.length} = 7
\]
(\texttt{length\_seven}).
The count is two components of multiplicity greater than one and five of multiplicity one, written $2+5$ in \texttt{ComponentsOfKodairaNeron.lean}.
For even $a$, \texttt{BlowupEvenA.lean} gives
\[
\texttt{EvenAComponent.all.length} = 9
\]
(\texttt{length\_nine}, \texttt{even\_a\_nine\_components}), written $4+5$.

\texttt{kodaira\_neron\_odd\_a\_open} and \texttt{kodaira\_neron\_even\_a\_open} are the propositions that identify these counts with the opaque constant \texttt{NeronSpecialFiberComponents}.
They are hypotheses, not theorems.
Adjacency of the components, normal crossings, minimality of the model, and the identification of the N\'eron model with the smooth locus are not proved.

\section{The rewired Pacetti path}
\texttt{pacetti\_2\_8\_open} is a live proposition.
On a coprime curve, it says that the sum of the sealed tame and Swan conductors at $\pi$ lies in $\{10,12\}$.
\texttt{pacetti\_of\_kodaira\_neron} proves that sentence from the two Kodaira--N\'eron opens above.
The capstone calls this theorem and passes the resulting term to the level-lowering axiom.
It calls no other derivation of the exponent.
\texttt{pacetti\_of\_components} still derives the same sentence from the separate open \texttt{components\_n\_plus\_}\allowbreak\texttt{five\_open}.
The capstone does not assume that open and does not call that theorem.

The conductor shape $N=32|c|^2$ passed to Ribet's axiom is the predicate \texttt{FreyConductorDivides}.
It is not a computation of the conductor.

\section{The Galois skeleton}
\texttt{DarmonMerel/8\_Galois4413.lean} contains definitions and theorem statements only.
It declares no \texttt{axiom} and uses no \texttt{sorry}.

\texttt{structure RhoE13} is stated for a \texttt{FreyCurve4413}.

The representation space is \texttt{RhoE13Space.unresolved}.
The residual is \texttt{rho\_E\_p}.
That residual is \texttt{GaloisRepresentation.unresolved}.
There is no formalization of $\mathrm{Gal}(\overline{\mathbb Q}/\mathbb Q(i))$ and none of $\mathrm{GL}_2(\mathbb F_{13})$.
\texttt{RhoE13.ofFrey} builds that unresolved package.
It does not produce a term of \texttt{RhoE13Open}.
\texttt{rhoE13\_irreducible\_open} is \texttt{RhoE13Open}.

\texttt{structure qcurve\_extension\_169} stores the conjugate of the linear coefficient and the equality \texttt{velu\_target\_eq\_twist\_neg\_two}.
Both sides of that equality are the coefficient pair $(-8a,\, 8a^2+8b^2\sqrt{-1})$.
The equality is the coefficient form of the $2$-isogeny from the Galois conjugate to the twist by $-2$.
It is not a class in a cohomology group and not a map on a Tate module.
\texttt{qcurve\_extension\_169.ofFrey} fills those fields from the proved coefficient theorems.
It does not prove \texttt{qcurve\_extension\_169\_open}.
\texttt{qcurve\_extension\_169\_cocycle\_open} is that proposition.

\texttt{frey\_qcurve\_modular\_on\_RhoE13} states what \texttt{frey\_qcurve\_modular} gives.
The conclusion is \texttt{Modular E}, the residual of \texttt{RhoE13.ofFrey} left unresolved, and a nonempty coefficient extension.
\texttt{ribet\_level\_lowering\_on\_RhoE13} states what \texttt{ribet\_level\_lowering\_4413} gives.
From the conductor shape, a term of \texttt{RhoE13Open}, and \texttt{pacetti\_2\_8\_open}, the conclusion is a nonempty \texttt{RibetLevel}, the residual isomorphism \texttt{GkResidualIso}, and the same unresolved residual.
Neither statement constructs a representation.

\section{Open inventory}
Five propositions remain open. None is an \texttt{axiom} and none is \texttt{sorry}.

\begin{enumerate}
\item \texttt{kodaira\_neron\_odd\_a\_open}, as above. Hypothesis of the capstone.
\item \texttt{kodaira\_neron\_even\_a\_open}, as above. Hypothesis of the capstone.
\item \texttt{no\_newform\_32\_64\_open}, defined in \texttt{DarmonMerel/3\_RibetLowering4413.lean} immediately after \texttt{structure RibetLevel}, and before \texttt{GkResidualIso}.
The capstone applies it once, to the package returned by the level-lowering axiom.
It is the missing \texttt{False} from Darmon--Merel Lemma~8.2. The Mazur products at the newforms $32.2.a.a$ and $64.2.a.a$ are $0$, and $(T_5-6)e\not\equiv 0\pmod{13}$ at both levels, so those two calculations do not yet contradict each other.
\item \texttt{RhoE13Open}, an \texttt{opaque} proposition with no body, in the same file.
Nothing in the repository produces a term. The level-lowering axiom does not apply without one.
The carrier is \texttt{RhoE13}, described above.
\item \texttt{qcurve\_extension\_169\_open}, stated as the trace divisibility $13\mid \texttt{freyTrace3}-(-6)$ for a lift through $\mathrm{GL}_2(\mathbb F_{169})$.
\texttt{no\_newform\_32\_64\_of\_qcurve} derives item~3 from this sentence.
The capstone does not assume it.
The structure \texttt{qcurve\_extension\_169} does not prove the divisibility.
\texttt{thirteen\_not\_dvd\_two}, the integrality $13\nmid 2$, is a theorem in file~3 and is used by the isogeny input of that derivation.
\end{enumerate}

\section{The two axioms}
\texttt{frey\_qcurve\_modular} produces \texttt{Modular E}.

\texttt{ribet\_level\_lowering\_4413} is a structure \texttt{RibetLowering}.

Its level field returns a \texttt{RibetLevel} from the conductor shape.
The hypotheses of that field are a term of \texttt{RhoE13Open} and \texttt{pacetti\_2\_8\_open}.
Its other field is the residual isomorphism \texttt{GkResidualIso}.
The disjunction $n=32\lor n=64$ is part of that axiom. It is not deduced from the component counts.
The skeleton does not add an axiom.

\section{Closing}
The library formalizes the Darmon--Merel method for the signature $a^4+b^4=c^{13}$, as far as the definitions, the component counts, the rewired exponent, the Galois skeleton, and the conditional capstone.
It is not a proof of $a^x+b^y=c^z$ for general exponents, and it is not a proof of $a^4+b^4=c^{13}$.
Whether that formalization is published is the editors' decision.
A DOI for the repository is a separate question: the default branch \texttt{main} is a different history, and this branch still carries the two axioms and the five opens above.
The only manuscript in this branch is this file.
\end{document}
